% Generated draft: requires mathematical and source review.
\documentclass[11pt]{article}
\usepackage[margin=25mm]{geometry}
\usepackage{fontspec}
\setmainfont{Latin Modern Roman}
\usepackage{amsmath,amssymb,mathtools,unicode-math}
\setmathfont{Latin Modern Math}
\usepackage{enumitem,needspace,xurl,hyperref}
\setlength{\parindent}{0pt}
\setlength{\parskip}{5pt}
\setlength{\emergencystretch}{3em}
\newcommand{\N}{\mathbb{N}}
\newcommand{\Z}{\mathbb{Z}}
\newcommand{\Q}{\mathbb{Q}}
\newcommand{\R}{\mathbb{R}}
\newcommand{\C}{\mathbb{C}}
\newcommand{\F}{\mathbb{F}}
\newcommand{\A}{\mathbb{A}}
\newcommand{\PP}{\mathbb{P}}
\newcommand{\E}{\mathbb{E}}
\newcommand{\eps}{\varepsilon}
\newcommand{\dd}{\,\mathrm{d}}
\begin{document}
\section*{UnsolvedMath ID: 10900008}
\textbf{Source problem code:} catalogue entry
\subsection{Definitions and mathematical statement}
Let \(V\) be a finite-dimensional real vector space and let \(\beta:V\times V\to\R\) be a nondegenerate bilinear form, not assumed symmetric or skew-symmetric. Set
\[
 G=\operatorname{Isom}(\beta)
 =\{A\in\operatorname{GL}(V):\beta(Au,Av)=\beta(u,v)\ \forall u,v\in V\}.
\]
In a basis with form matrix \(B\), this is \(\{A:A^{\mathsf T}BA=B\}\).

\emph{Problem (Cooper).} Classify the subgroups \(H\subseteq\operatorname{GL}(V)\) for which there is a sequence \(A_n\in\operatorname{GL}(V)\) such that
\[
 A_nGA_n^{-1}\longrightarrow H
\]
in the Hausdorff-limit sense intended by the source. In the usual closed-subgroup, locally Hausdorff (Chabauty) interpretation, the limit \(H\) is closed and the convergence can be stated as follows:
\begin{enumerate}
\item For every \(h\in H\), there exist \(h_n\in A_nGA_n^{-1}\) with \(h_n\to h\).
\item If \(n_k\to\infty\) and \(h_{n_k}\in A_{n_k}GA_{n_k}^{-1}\) converge in \(\operatorname{GL}(V)\) to \(h\), then \(h\in H\).
\end{enumerate}
These conditions specify a standard interpretation, not an unprovided choice of global Hausdorff metric on an unbounded space. The source says the classification is known when \(\beta\) is nondegenerate symmetric or nondegenerate skew-symmetric; its question retains arbitrary nondegenerate real bilinear forms. The excerpt gives no reference for those known cases and does not explicitly define its Hausdorff topology.
\subsection{Short English statement}
For any nondegenerate real bilinear form, classify all limits of conjugates of its isometry group inside \(\operatorname{GL}(V)\). The symmetric and skew-symmetric cases are reported known; the general form and the exact noncompact Hausdorff-limit convention must be retained.
\subsection{Sources}
\begin{itemize}
\item Local UnsolvedMath registry, permanent ID 10900008, source code catalogue entry; attribution Cooper, subsection 2.3, ``Hausdorff limits of conjugates of an isometry group.''
\item Delp, Hoffoss and Manning, \emph{Problems In Groups, Geometry, and Three-Manifolds} (2015), subsection 2.3, PDF page 2. Supplied URLs: \url{https://arxiv.org/pdf/1512.04620#page=2} and \url{https://arxiv.org/abs/1512.04620}; recorded access 2026-07-29, extraction source-tex checked against rendered PDF text. The source reports the symmetric and skew-symmetric cases known.
\item The research summary mentions Breuillard's ``Diophantine geometry and uniform growth'' and related work, approximate subgroups, and algebraic groups over local fields as possible context, without precise results resolving this classification or full bibliographic details. Its mention of semisimple groups does not impose semisimplicity on the problem.
\item Supplied status open and OPEN-TRIAGE. The summary reports no direct published classification found; background records that the 2015 source presents the problem as open and no complete modern primary-literature audit is recorded. Status and redistribution terms for the publicly accessible author-submitted TeX and PDF remain NEEDS\_REVIEW.
\end{itemize}
\end{document}
